\documentclass[11pt]{article}

\usepackage[margin=1.2in]{geometry}
\usepackage{amsmath,amssymb}
\usepackage{microtype}
\usepackage{booktabs}
\usepackage{xcolor}
\usepackage{mdframed}
\usepackage[colorlinks=true,linkcolor=blue!50!black,urlcolor=blue!50!black]{hyperref}

\newmdenv[backgroundcolor=blue!4,linecolor=blue!40,linewidth=1pt,
  skipabove=8pt,skipbelow=8pt]{keyidea}

\title{An Improved Lower Bound for the $3x+1$ Problem:\\
       \Large A Friendly Guide}
\author{Companion notes to ``$\pi_a(x) \ge x^{0.895}$''\thanks{Written
in collaboration with Claude (Anthropic). These notes explain the ideas;
the paper contains the precise statements and the certification
procedure. The Collatz conjecture itself remains open.}}
\date{June 12, 2026}

\begin{document}
\maketitle

\section{The question}

Nobody can prove that \emph{every} number reaches 1 under the Collatz
rule (halve if even, triple-and-add-one if odd). So mathematicians ask a
humbler question with a provable answer:

\begin{center}
\emph{Out of the first $x$ numbers, how many can we PROVE reach 1?}
\end{center}

Call that count $\pi_1(x)$. If the conjecture is true, $\pi_1(x) = x$.
What can actually be proved has crept up over five decades, always in
the form ``at least $x^{\beta}$ of them'':
\begin{center}
\begin{tabular}{lll}
\toprule
year & authors & exponent $\beta$ \\
\midrule
1978 & Crandall & small positive \\
1989 & Krasikov & $0.43$ \\
1998 & Wirsching & $0.48$ \\
1995 & Applegate--Lagarias & $0.81$ \\
2003 & Krasikov--Lagarias & $0.84$ \\
2026 & this work & $\mathbf{0.895}$ \\
\bottomrule
\end{tabular}
\end{center}
To calibrate: $x^{0.895}$ of a trillion is about 52 billion, versus 13
billion for $x^{0.84}$ --- and the dream, equivalent to ``almost
density one,'' would be $x^{0.999\ldots}$. The 2003 record stood for
twenty-three years. The improvement here required, as the paper's
acknowledgments put it, \emph{no new mathematics --- only twenty-three
years of Moore's law and one change of algorithm}. These notes explain
the method from scratch, because it is unusually pretty.

\section{Growing the tree backwards}

Instead of asking ``does $n$ reach 1?'', turn the dynamics around:
start at 1 and ask \emph{which numbers arrive here?} Every number $n$
has one guaranteed ancestor-in-reverse, namely $2n$ (because $2n$
halves to $n$). Some numbers have a second one: if $2n - 1$ is
divisible by 3, then $b = (2n-1)/3$ is an odd number whose
triple-plus-one, halved, is exactly $n$.

So the numbers that provably reach 1 form a \textbf{tree}: from each
node, one branch always continues (doubling), and a second branch
sprouts whenever a divisibility condition holds. Every node of the tree
is a certified success --- no conjecture needed. If the tree branches
often enough, it contains exponentially many numbers, and counting them
gives a lower bound for $\pi_1(x)$.

There is a tension, though. Doubling makes tree members \emph{bigger},
so a tree with many members might spend them on numbers far beyond
$x$. What matters is the count of members \emph{below} $x$: roughly,
(how fast the tree branches) measured against (how fast its members
grow). That ratio is what becomes the exponent $\beta$.

\begin{keyidea}
\textbf{Who gets the second child?} The condition ``$2n-1$ divisible by
3'' depends only on the remainder of $n$ when divided by 3. And the
remainder of the child depends on the remainder of the parent when
divided by 9, and so on: remainders modulo $3, 9, 27, \ldots, 3^k$ are
the natural coordinates of the tree. Tracking finer remainders (larger
$k$) gives a more detailed map of where branching is plentiful --- and
provably never a worse exponent. The entire fifty-year history in the
table above is, in essence, the same tree examined at finer and finer
resolution.
\end{keyidea}

\section{From tree to exponent: the certificate}

Here is the 2003 framework, de-jargonized. Assign to every remainder
class $m$ (mod $3^k$) a positive weight $c_m$ --- think of it as a
prediction of how fertile that class's subtrees are. Krasikov and
Lagarias wrote down a system of inequalities --- one per class, each
saying ``this class's weight is justified by the weights of the classes
its children land in, discounted by a growth rate $\lambda$'' --- and
proved:

\begin{center}
\emph{If weights satisfying the system exist for some rate $\lambda$,
then $\pi_1(x) \ge x^{\log_2 \lambda}$.}
\end{center}

The weights are called a \textbf{certificate}: a finite list of numbers
that, once checked against the finitely many inequalities, proves an
infinite statement. Bigger $\lambda$ = better exponent; finer classes
(bigger $k$) admit bigger $\lambda$. In 2003 the state of the art in
solving such systems (linear programming) reached $k = 11$ --- 59{,}049
classes --- and $\lambda = 1.792$, i.e.\ the exponent $0.84$.

\section{The unlock: certificates that find themselves}

The new observation is about the \emph{shape} of the system. For a
fixed rate $\lambda$, the inequalities say precisely
\[
c \;\le\; F(c),
\]
where $F$ takes the whole list of weights and produces a new list, each
entry built from other entries by adding, scaling by positive
constants, and taking minimums. Such an $F$ has two golden properties:
it is \emph{monotone} (raise inputs, outputs never drop) and
\emph{scale-invariant} (double all inputs, outputs double).

For maps like that there is a classical trick --- ironically associated
with the same Lothar Collatz, in his day job as a numerical analyst.
Start from any guess, apply $F$, rescale, repeat (readers who know how
search engines rank pages have seen this dance: it is power iteration).
The guesses settle toward an equilibrium, and the punchline is:

\begin{keyidea}
\textbf{Any non-shrinking guess is a proof.} You do not need to solve
the system, trust the iteration, or invoke any theory: if at any moment
your list satisfies $c \le F(c)$ entry-by-entry --- ``one more
application of the rules doesn't shrink any weight'' --- then that very
list is a valid certificate, checkable by anyone with a multiplication
routine. The iteration is just a way of \emph{finding} such a list; it
plays no role in the proof.
\end{keyidea}

One sweep of $F$ costs time proportional to the number of classes. That
replaces the 2002 bottleneck (general-purpose linear programming) and
carries the computation from $59$ thousand classes to $43$ million
($k = 17$) on a laptop, with two further accelerations: each level's
certificate is recycled as the starting guess for the next (the 2003
paper itself proves certificates lift between levels), and at the
largest levels one simply aims at a target rate directly --- a hit is a
hit, no search required.

\section{Why you can trust a 43-million-line proof}

Floating-point arithmetic found the certificates; floating-point
arithmetic is famously not to be trusted. So nothing in the final claim
rests on it:

\begin{enumerate}
\item \textbf{Rational rates.} The claimed rate is chosen as
$\lambda = 2^{p/q}$ for whole numbers $p, q$ --- so the exponent is
\emph{exactly} $p/q$ (for the record: $179/200 = 0.895$). The system's
coefficients involve irrational quantities; each is replaced by a
nearby fraction that is provably \emph{below} it, the proof being a
single comparison of huge whole numbers (is $u^q \cdot 2^{2p} \le
v^q \cdot 3^p$?). Using smaller coefficients only makes the
inequalities harder to satisfy --- errors can only hurt us, never help
--- so soundness leans one way.
\item \textbf{Exact verification.} The certificate is converted to
whole numbers and every one of the $43{,}046{,}721$ inequalities is
checked in exact integer arithmetic. No rounding exists anywhere in the
verified statement.
\item \textbf{Two independent referees.} The checker was implemented
twice --- the second time from a clean slate, re-deriving the
inequality system from the published paper, using different number
representations and different bounding tricks, sharing not a line of
code. Both pass every level. And the pair was tested for blindness:
corrupt a single entry of a certificate and both checkers catch
exactly that violation.
\item \textbf{Validation against history.} Run at the 2003 levels, the
machinery reproduces every entry of Krasikov--Lagarias' published table
to all seven printed decimals, including their record $1.7922310$.
\end{enumerate}

The mathematical heavy lifting --- the theorem that a certificate
yields the bound at all --- is exactly the published, peer-reviewed
2003 theorem, used verbatim. This work adds no asymptotic analysis; it
adds certificates the 2003 authors could not compute, plus the means to
check them beyond reasonable doubt.

\section{The extra mile: a fragment proved by a machine, end to end}

One can go one step further in trustworthiness: have a proof assistant
verify not just the certificate but the \emph{theorem} too. The
companion repository contains a complete formalization, in Lean 4, of a
simplified variant of the whole chain --- the backward tree, the
inequalities (proved from the dynamics, including the amusing fact that
the tree can contain no loops because the orbit of 8 ends in the
$1 \to 2 \to 1$ cycle), an 81-class certificate checked by Lean's own
kernel, and the resulting bound
\[
\pi_1(x) \;\geq\; \text{const} \cdot x^{0.4543},
\]
with literally every logical step machine-checked. That exponent is
deliberately modest --- the simplification trades strength for
formalizability --- but it already exceeds every bound published before
1995, and it is, to our knowledge, the first machine-verified density
bound for this problem in history. Scaling the formalization up toward
$0.895$ is bookkeeping rather than new ideas.

\section{What it would take to finish}

Within this method, everything hinges on one open question: as the
class resolution $k \to \infty$, does the best rate $\lambda_k$ climb
all the way to 2? If yes, the exponent climbs to 1 and \emph{almost
all} numbers provably reach 1 in the strongest counting sense. The
data is consistent with it --- the certified rates march
$1.79, 1.81, 1.82, 1.83, 1.84, 1.85, 1.86$ with slowly shrinking steps
--- but each new level costs triple the memory and compute, and the
limit question itself looks like genuine mathematics, not computation.
Meanwhile the companion paper explains why no method that examines
numbers through any bounded lens can settle the full conjecture: the
last mile belongs to ideas not yet had.

\section{Summary in four sentences}

Count the numbers provably reaching 1 by growing the tree of ancestors
of 1, organized by remainders modulo $3^k$. A finite list of weights
satisfying finitely many inequalities certifies an infinite bound
$\pi_1(x) \ge x^{\log_2 \lambda}$ --- and since the system is monotone
and scale-invariant, certificates can be \emph{found} by simple
iteration and \emph{verified} by exact arithmetic, no linear
programming and no floating point in the final claim. Pushing from
59 thousand to 43 million classes lifts the twenty-three-year-old
record from $x^{0.84}$ to $x^{0.895}$, checked by two independent
verifiers and reproducing the historical table exactly. A simplified
fragment is additionally machine-verified end to end in Lean --- the
first formally verified density bound for the problem --- at
$x^{0.4543}$.

\end{document}
